\documentclass[UTF8,zihao=-4]{ctexart}
\usepackage[a4paper,margin=2.2cm]{geometry}
\usepackage{xcolor}
\usepackage{hyperref}
\usepackage{enumitem}
\hypersetup{colorlinks=true,linkcolor=teal,urlcolor=teal}
\setlist{itemsep=0.25em,topsep=0.35em}
\title{\bfseries 点子发芽日报 2026-05-25}
\author{点子发芽}
\date{2026-05-25}
\begin{document}
\maketitle
\section*{今日总结}
今天没有新的网页输入，所以日报不应该硬写成新闻汇总。更有价值的动作是检查“新知孵化工作流”这个主线是否仍然成立：如果以后只通过网页输入关键词、上下文和权重，报告就必须足够短、足够清楚、足够像人写的。
\section*{今日输入}
今日没有新的网页输入。
\section*{今日新知}
\subsection*{新知孵化工作流}
\paragraph{简介} 新知孵化工作流是这个项目的核心节点：它要把白天遇到的新信息变成晚上能读、明天能做的小判断。今天没有新输入，反而暴露出一个重要边界：系统不能为了显得勤奋而编一堆复盘内容，没材料时就应该短一点。
\paragraph{最近有什么相关新闻} 这次外部搜索主要找到个人知识管理、Markdown 笔记和每日复盘类普通资料，适合作为背景，不适合作为新的事实结论 [S1] [S6] [S11]。因此今天不写所谓新进展，只保留一个判断：报告结构越短，越可能被你长期读下去。
\paragraph{最小下一步} 明天从网页提交 3 个真实关键词，然后只检查一件事：报告里的每个关键词是否都能用“简介、相关新闻、下一步”讲清楚。
\section*{与旧知识的链接}
\begin{itemize}[leftmargin=2em]
\item 这个节点直接连接 README、automation 提示词和日报生成器，是当前项目最重要的长期主题。
\item 它也连接“每日小推进”点子：日报最后应该推动一个很小的动作，而不是生产一堆看起来完整的废话。
\end{itemize}
\section*{今日发芽点子}
\begin{itemize}[leftmargin=2em]
\item 报告瘦身实验：连续三天只保留简版结构，观察你是否更愿意读完。
\end{itemize}
\section*{参考搜索内容}
\begin{itemize}[leftmargin=2em]
\item \href{https://www.zhihu.com/tardis/zm/art/1955849427556824396}{[S1] 普通线索：AI时代的个人知识管理最佳实践：从笔记仓库到认知操作系统}
\item \href{https://github.com/flyhunterl/flymd}{[S6] 项目主页/机构主页：GitHub - flyhunterl/flymd: 高性能Markdown笔记工具}
\item \href{https://juejin.cn/post/7571113889277194252}{[S11] 普通线索：提升个人知识整理与复盘能力的实用方法}
\end{itemize}
\end{document}
